\chapter{The Cache Hierarchy, False Sharing, and Cache-Conscious Layout}

For most real workloads the bottleneck is not computing but \emph{feeding} the computation with data, and the cache hierarchy is the apparatus for doing so. A single miss to main memory costs $\sim$300 instructions; whether your data is laid out so the hardware can predict and prefetch it is the difference between fast and unusably slow. This chapter explains the cache --- lines, levels, associativity, prefetching, coherence --- and turns it into layout disciplines: locality, hot/cold splitting, and the elimination of false sharing.

\section*{Chapter Roadmap}
\begin{itemize}
\item 87.1 Why the Cache Exists: The Memory Wall
\item 87.2 Cache Lines: The Unit of Transfer
\item 87.3 The Levels: L1, L2, L3, and Inclusivity
\item 87.4 Associativity and Conflict Misses
\item 87.5 Spatial and Temporal Locality
\item 87.6 Hardware Prefetching
\item 87.7 Cache Coherence and the Cost of Sharing
\item 87.8 False Sharing and Interference Sizes
\item 87.9 Cache-Conscious Data Layout
\end{itemize}

\section{Why the Cache Exists: The Memory Wall}
A core executes $\sim$4 instructions per $\sim$0.3\,ns cycle, but DRAM takes $\sim$100\,ns ($\sim$1000+ instructions). Caches are small fast SRAM holding recently/soon-used data near the core.

\textbf{Why this matters.} Fast memory is small; large memory is slow. You cannot make DRAM faster, but you can change your hit rate by layout and traversal --- and hit rate governs memory-bound performance.

\section{Cache Lines: The Unit of Transfer}
The cache moves 64-byte lines, not bytes. Touching one byte fetches its whole line; adjacent data comes for free.

\textbf{Why this matters / cost model.} The 64-byte line is the key number: a sub-line struct accessed fully costs one miss; a contiguous array amortises one miss over many elements; a scattered linked list pays a miss per node. The line is also the unit of coherence (hence false sharing). Most layout rules follow from the 64-byte granularity.

\section{The Levels: L1, L2, L3, and Inclusivity}
\begin{itemize}
\item L1 (split I/D): 32--64\,KB each, $\sim$4 cycles, per core.
\item L2: 256\,KB--2\,MB, $\sim$12 cycles, per core.
\item L3 (LLC): 8--64\,MB, $\sim$40 cycles, shared.
\item DRAM: GBs, $\sim$100\,ns, all cores.
\end{itemize}

\textbf{Why this matters.} The order-of-magnitude steps make working-set size matter: a loop fitting in L1 runs an order faster than one spilling to L2, another faster than DRAM. Blocking/tiling keeps the active set in L1/L2. Know your structure's size relative to the levels.

\section{Associativity and Conflict Misses}
An N-way set-associative cache maps each address to a set of N slots; more than N hot lines in one set evict each other (a conflict miss) despite free capacity elsewhere.

\textbf{Why this matters / cost model.} This explains the power-of-two stride pathology: large power-of-two strides map many addresses to one set, thrashing it. Fix by padding array dimensions to avoid such strides. A profiler-only bug; awareness lets you recognise it.

\section{Spatial and Temporal Locality}
Spatial: nearby addresses accessed soon (contiguous/sequential access). Temporal: the same address reused soon (loop blocking, small hot data).

\textbf{Why this matters.} These explain the recurring advice: \texttt{vector} over \texttt{list} (spatial), SoA for partial-field access (Chapter 90), and loop blocking (temporal). Cache-conscious design is the pursuit of both.

\section{Hardware Prefetching}
Prefetchers detect sequential/strided patterns and fetch lines ahead of use, hiding DRAM latency.

\begin{lstlisting}[language=C++, caption={Sequential access is prefetched; random access is not. Min standard: C++11.}]
long sum = 0;  for (size_t i=0;i<n;++i) sum += a[i];            // unit stride: prefetched
long sum2 = 0; for (size_t i=0;i<n;++i) sum2 += a[index[i]];    // random: a miss per access
\end{lstlisting}

\textbf{Why this matters / cost model.} A sequential 1\,GB scan is near bandwidth-limited; random 1\,GB access pays a $\sim$100\,ns miss each --- 50$\times$ slower for the same data. This is why pointer-chasing is slow: each \texttt{->next} is a data-dependent load the prefetcher cannot predict. Make access sequential first.

\section{Cache Coherence and the Cost of Sharing}
Coherence (MESI-family) keeps cores consistent: before writing a line, a core gains exclusive ownership, invalidating other copies; another core's read then re-fetches.

\textbf{Why this matters / cost model.} Writes to shared lines cost a coherence round-trip (tens to hundreds of cycles); a line written by one core and read by another ping-pongs. This is the reality beneath the memory model: an atomic RMW must win exclusive ownership. Minimise shared mutable lines (thread-per-core, Chapter 96).

\section{False Sharing and Interference Sizes}
Two cores writing logically independent variables on one 64-byte line: each write invalidates the other's line though no data is shared.

\begin{lstlisting}[language=C++, caption={Padding to destructive interference size eliminates false sharing. Min standard: C++17.}]
struct Bad { std::atomic<long> a; std::atomic<long> b; };       // both within 64 bytes
struct Good {
    alignas(std::hardware_destructive_interference_size) std::atomic<long> a;
    alignas(std::hardware_destructive_interference_size) std::atomic<long> b;
};
\end{lstlisting}

C++17 \texttt{<new>} provides \texttt{hardware\_destructive\_interference\_size} (offset to separate lines --- avoid false sharing) and \texttt{hardware\_constructive\_interference\_size} (size to keep on one line --- promote locality).

\textbf{Why this matters / cost model.} False sharing can make a parallel program slower than serial: N per-thread counters on one line ping-pong on every increment. Invisible in source; only \texttt{perf c2c} or a layout audit reveals it. Pad each hot per-thread datum to its own line.

\section{Cache-Conscious Data Layout}
\begin{lstlisting}[language=C++, caption={Hot/cold splitting keeps the hot working set dense. Min standard: C++11.}]
struct OrderBad  { double price; int qty; char description[200]; };   // cold bloats the line
struct OrderGood { double price; int qty; const char* description; }; // hot part < 1 line
\end{lstlisting}

Layout rules: contiguous over linked; hot/cold splitting; structure-of-arrays for cross-object field access; pack hot data within a line and pad shared-mutable data across lines; size working sets to the cache via blocking.

\textbf{The discipline.} Memory-bound performance is won at design time by layout. Ask of every hot structure: contiguous? does a line carry only used bytes? are per-thread writes on separate lines? does the working set fit in cache? These deliver the largest gains in this volume (2--50$\times$).
